\documentclass[11pt]{article}
\usepackage[margin=1in]{geometry}
\usepackage{amsmath, amssymb}
\usepackage{booktabs}
\usepackage{caption}
\usepackage{multirow}
\usepackage{xcolor}
\usepackage[colorlinks=true, linkcolor=blue, citecolor=blue]{hyperref}

\newcommand{\bs}{\boldsymbol{s}}
\newcommand{\R}{\mathbb{R}}
\newcommand{\E}{\mathbb{E}}
\newcommand{\one}{\mathbb{1}}
\newcommand{\Kmrts}{K_{\mathrm{MRTS}}}

\title{Non-stationary simulation study:\\
what a mean-only MRTS augmentation can and cannot buy}
\author{}
\date{July 2026 --- interim report v3 (settings 1--12, the
threshold-vs-symmetry disambiguation, and the $K \ge 30$ extension)}

\begin{document}
\maketitle

\begin{abstract}
We evaluate the MRTS covariate extension of the Morris et al.\ (2017) skew-$t$
space--time model on twelve simulated data-generating processes (DGPs) that place
non-stationarity in different moments: the mean (settings 1, 2, 11, 12), the spatial
dependence (settings 3--8), the tails (setting 9), and the skewness (setting 10).
Three findings.
(i)~MRTS improves predictive scores \emph{only} where the truth carries
time-averaged mean structure, and this is provable in advance: an oracle
calculation shows the relative Brier score has a floor $R^\star \approx 1$ for
nine of eleven settings, so their observed nulls are properties of the design,
not evidence about MRTS.
(ii)~Where mean structure exists, the energy- and variogram-score gains are
carried \emph{entirely} by the two thresholded symmetric-$t$ methods. A
disambiguating pair of methods (symmetric-$t$ \emph{without} threshold,
$n=10$) attributes this cleanly to the \emph{threshold}, not the symmetry ---
and the absolute scores reframe it: POT censoring without a mean basis damages
the predictive field so badly (baseline energy $2.6\times$--$16\times$ worse)
that MRTS is a partial \emph{rescue}, not a synergy; the unthresholded methods
never needed rescuing and remain better in absolute terms at every $K$.
(iii)~The Brier score is close to blind to mean recovery and can even
deteriorate with $K$: in the strong-mean setting~11 a fit that recovers every
parameter of the DGP at $K=25$ has a \emph{worse} Brier score than the badly
misspecified $K=0$ baseline. The rich-mean setting~12, run out to $K=50$,
resolves this: the deterioration is confined to intermediate $K$ (where the
mean basis is still too coarse), and once the basis is fine enough
($K \ge 30$--$40$) the Brier score improves by 7--12\%, consistent with the
extrapolation-error mechanism. The discriminating check passes: the
simple-surface setting~11, extended to the same $K = 50$, never turns
positive --- its Brier ratio only decays back to parity ($1.09 \to 1.00$) ---
so the reversal requires mean structure that only a fine basis can resolve,
not merely a large $K$.
\end{abstract}

\section{Design}

\paragraph{Data-generating process.}
All settings share the skew-$t$-1 core on $n_s=144$ sites drawn uniformly on
$[0,10]^2$:
\begin{equation}
Y_t(\bs) \;=\; 10 \;+\; \lambda\,\sigma_t\,|z_t| \;+\; \sigma_t\, v_t(\bs)
\;+\; [\text{structure}]_k,
\qquad t = 1,\dots,50,
\label{eq:core}
\end{equation}
with $\lambda = 3$, $z_t \stackrel{iid}{\sim} N(0,1)$ half-normal,
$v_t \sim N(0, C)$ exponential correlation ($\rho=1$, $\gamma=0.9$),
$\tau_t \sim \mathrm{Ga}(3/2, 4)$, $\sigma_t = \tau_t^{-1/2}$ (marginal $t_3$),
i.i.d.\ in time. The settings differ only in the structure term---or, for the
\emph{replace-$C$} settings, in the error correlation itself
(Table~\ref{tab:dgp}). Table~\ref{tab:anatomy} gives the resulting anatomy of
each DGP: the marginal law of $Y_t(\bs)$, the spatial dependence of the error
process, and the time-averaged true mean surface $\bar\mu(\bs)$ (the target of
any pooled-$\beta$ mean basis).

\begin{table}[ht]
\centering\small
\caption{The twelve fitted DGP settings and the number of simulation replicates
$n$ (independent datasets) fitted per setting. Routes: \emph{mean} = added to
the fixed mean; \emph{additive} = mean-zero random field added to the draw;
\emph{replace $C$} = non-stationary correlation handed to the generator (unit
diagonal, so the \emph{marginal} skew-$t$ law is unchanged --- verified by KS
tests, $p \in [0.075, 0.864]$); \emph{skew term} = spatially varying $\lambda$.}
\label{tab:dgp}
\begin{tabular}{llllc}
\toprule
ID & \texttt{surf\_type} & non-stationarity & route & $n$ \\
\midrule
1  & invariant         & mean, 2 cosine bumps, $\mathrm{sd}=3.7$   & mean       & 10 \\
2  & varying           & mean, time-varying weights                & mean       & 10$^{a}$ \\
3  & ns\_dependence    & dependence, low-rank cosine RE            & additive   & 10 \\
4  & ps\_cov           & dependence, Paciorek--Schervish           & replace $C$& 3  \\
5  & ps\_add           & dependence, same PS field                 & additive   & 3  \\
6  & covdep\_cov       & dependence, $\rho(\bs)=e^{0.7 f_1(\bs)}$  & replace $C$& 3  \\
7  & covdep\_add       & dependence, same field                    & additive   & 3  \\
8  & fuentes\_cov      & dependence, 4-regime mixture              & replace $C$& 3  \\
9  & gh\_marginal      & tails, Tukey $g$-and-$h$, $g(\bs), h(\bs)$& additive   & 3  \\
10 & lambda\_varying   & skewness, $\lambda(\bs) = 3 + 3 f_2(\bs)$ & skew term  & 3  \\
11 & invariant\_strong & mean, $3\times$ setting 1, $\mathrm{sd}=11$ & mean     & 3  \\
12 & mrts\_multibump   & mean, 6 Gaussian bumps, $\mathrm{sd}=11$    & mean     & 3  \\
\bottomrule
\multicolumn{5}{l}{\footnotesize $^{a}$ setting 2 was fitted at $K=0$ only (no MRTS sweep); it is excluded from the score tables.}
\end{tabular}
\end{table}

\begin{table}[ht]
\centering\footnotesize
\caption{Anatomy of the twelve DGPs: marginal law of $Y_t(\bs)$, spatial
dependence of the error process, and time-averaged true mean $\bar\mu(\bs)$
(its spatial sd in parentheses; per-dataset values for random means).
$C_{\mathrm{exp}}(h) = 0.9\, e^{-h}$ is the shared stationary exponential
correlation (Mat\'ern $\nu = \tfrac12$, $\rho = 1$, unit diagonal, i.e.\ a 10\%
nugget); ``skew-$t$'' is the $t_3$ core of \eqref{eq:core} with $\lambda = 3$;
$\ast$ denotes convolution with an independent additive field
($\sigma_{\mathrm{add}} = 3$); $f_1, f_2$ are the two cosine bumps.}
\label{tab:anatomy}
\begin{tabular}{clp{3.3cm}p{4.4cm}p{3.9cm}}
\toprule
ID & \texttt{surf\_type} & marginal law & spatial dependence & mean structure $\bar\mu(\bs)$ \\
\midrule
1  & invariant         & skew-$t$ & $C_{\mathrm{exp}}$, stationary & $10 + 5f_1 + 3f_2$ \hfill (3.7) \\
2  & varying           & skew-$t$ & $C_{\mathrm{exp}}$, stationary & $10 + w_{1t}f_1 + w_{2t}f_2$, $(w_{1t},w_{2t}) \sim N(0,\mathrm{diag}(25,9))$ per day; $\bar\mu \approx 10$ \hfill (0.4) \\
3  & ns\_dependence    & skew-$t$ $\ast$ Gaussian & $C_{\mathrm{exp}}$ $+$ low-rank cosine $\sum_j \tau_j^2 \phi_j(\bs)\phi_j(\bs')$; non-stationary, negative at long range & $10$ (RE mean-zero) \hfill (0.4) \\
4  & ps\_cov           & skew-$t$ (unchanged) & Paciorek--Schervish locally anisotropic kernel; scale, aspect and rotation each vary across the domain & $10$ exactly \hfill (0) \\
5  & ps\_add           & skew-$t$ $\ast$ Gaussian & $C_{\mathrm{exp}}$ $+$ $\sigma_{\mathrm{add}}^2\, C_{\mathrm{PS}}$, additive & $10$ (RE mean-zero) \hfill (0.4) \\
6  & covdep\_cov       & skew-$t$ (unchanged) & isotropic PS kernel with local range $\rho(\bs) = e^{0.7 f_1(\bs)} \in [0.5, 2.0]$ & $10$ exactly \hfill (0) \\
7  & covdep\_add       & skew-$t$ $\ast$ Gaussian & $C_{\mathrm{exp}}$ $+$ $\sigma_{\mathrm{add}}^2\, C_{\rho(\bs)}$, additive & $10$ (RE mean-zero) \hfill (0.4) \\
8  & fuentes\_cov      & skew-$t$ (unchanged) & Fuentes 4-regime weighted GP mixture, ranges $0.4$--$4$ & $10$ exactly \hfill (0) \\
9  & gh\_marginal      & skew-$t$ $\ast$ std.\ Tukey $g$-and-$h$; $g(\bs) = 0.8 f_1$, $h(\bs) = 0.12(1{+}f_2)$ & $C_{\mathrm{exp}}$ in both components (transform perturbs correlation only mildly) & $10$ (standardised mean-zero) \hfill (0.4) \\
10 & lambda\_varying   & skew-$t$ with $\lambda(\bs) = 3 + 3f_2 \in [0,6]$ & $C_{\mathrm{exp}}$, stationary & induced: $10 + \lambda(\bs)\,\bar z$, $\bar z = \frac{1}{T}\sum_t \sigma_t |z_t|$ \hfill (2.8) \\
11 & invariant\_strong & skew-$t$ & $C_{\mathrm{exp}}$, stationary & $10 + 15f_1 + 9f_2$ \hfill (11.0) \\
12 & mrts\_multibump   & skew-$t$ & $C_{\mathrm{exp}}$, stationary & $10\,+$ six alternating-sign Gaussian bumps (width 1.6) \hfill (11.0) \\
\bottomrule
\end{tabular}
\end{table}

\paragraph{Fitted methods and grid.}
Five baseline methods are fitted to every dataset:
(1) Gaussian; (2) skew-$t$, $K_{\text{knots}}=1$; (3) symmetric-$t$,
$K_{\text{knots}}=1$, POT threshold at $q(0.80)$; (4) skew-$t$,
$K_{\text{knots}}=5$; (5) symmetric-$t$, $K_{\text{knots}}=5$, threshold
$q(0.80)$. Each is run across MRTS design widths
$\Kmrts \in \{0, 3, 5, 8, 10, 12, 15, 20, 25\}$, where $\Kmrts = 0$ denotes
the baseline design $[1, s_1, s_2]$ and $\Kmrts = K$ replaces the design with
the $K$-column \texttt{autoFRK::mrts} basis built on the \emph{training} sites
only (column 1 = intercept, 2--3 linear, $4{:}K$ thin-plate-spline
eigenfunctions; $K=3$ spans the baseline). MCMC: 20{,}000 iterations, 10{,}000
burn-in. The random seed is shared across $\Kmrts$ within each
(method, dataset), so all $K$-comparisons are \emph{paired}. The grid is
being extended by $\Kmrts \in \{30, 40, 50\}$ setting by setting (complete
for settings 1, 4, 11, 12 at the time of writing; Table~\ref{tab:brierK}
shows exactly which cells exist).

\paragraph{Scoring.}
The last 44 of the 144 sites are held out entirely; all scores are computed
out-of-sample on those sites. With $\hat F$ the posterior predictive law:
\begin{align}
\mathrm{BS} &= \tfrac{1}{11}\sum_{p}\ \tfrac{1}{44\cdot 50}\sum_{i,t}
\big(\hat P(Y_t(\bs_i^*) > u_p) - \one\{y_t(\bs_i^*) > u_p\}\big)^2,
\quad p \in \{0.90, \dots, 0.995\},\\
\mathrm{ES} &= \tfrac{1}{50}\sum_t \Big[\E\,\lVert \boldsymbol X_t - \boldsymbol y_t\rVert
- \tfrac12 \E\,\lVert \boldsymbol X_t - \boldsymbol X_t'\rVert\Big],
\qquad \boldsymbol X_t, \boldsymbol X_t' \sim \hat F_t \ \text{(44-dim)},\\
\mathrm{VS} &= \tfrac{1}{50}\sum_t \sum_{i<j}
\big(|y_{it}-y_{jt}|^{1/2} - \E\,|X_{it}-X_{jt}|^{1/2}\big)^2 .
\end{align}
The Brier score is a per-cell \emph{marginal} score and cannot see spatial
dependence; the energy score is the multivariate CRPS (reduces to CRPS at
$d=1$) and is location-sensitive; the variogram score ($p=1/2$) targets
pairwise dependence. Every result below is the \emph{relative} score
$R_{m,K,d} = S_{m,K,d} / S_{m,0,d}$ against the same method's own $K=0$ fit,
averaged over datasets; $R<1$ means MRTS helps.

\section{Results}

Tables~\ref{tab:energy}--\ref{tab:vario} report, per setting and method, the
best (smallest) $\bar R_{m,K}$ over $K>0$, with the minimising $K$ in
parentheses. Asterisks mark $|\bar R - 1| > 2\,\widehat{\mathrm{SE}}$ (SE over
datasets; with $n=3$ this is indicative only). Table~\ref{tab:brierK} gives the
full setting~$\times$~$K$ profile of the relative Brier score, method by
method.

\begin{table}[ht]
\centering\small
\caption{\textbf{Energy score}, best over $K>0$ relative to the same method's
$K=0$ (minimising $K$ in parentheses). Values $<1$: MRTS helps; bold:
$\bar R < 0.75$.}
\label{tab:energy}
\setlength{\tabcolsep}{4pt}
\begin{tabular}{clcccccc}
\toprule
ID & non-stationarity & $n$ & Gaussian & Skew-$t$ K1 & Sym-$t$ K1 T & Skew-$t$ K5 & Sym-$t$ K5 T \\
\midrule
1 & mean (2 bumps, sd 3.7) & 10 & $0.994^{*}$ {\tiny(40)} & $0.980^{*}$ {\tiny(40)} & $\mathbf{0.579^{*}}$ {\tiny(20)} & $0.983^{*}$ {\tiny(50)} & $\mathbf{0.659^{*}}$ {\tiny(12)} \\
3 & dependence (low-rank) & 10 & $1.000$ {\tiny(5)} & $1.000$ {\tiny(3)} & $0.993$ {\tiny(5)} & $1.000$ {\tiny(15)} & $0.911$ {\tiny(3)} \\
4 & dependence (PS, repl.\ $C$) & 3 & $0.999^{*}$ {\tiny(5)} & $1.000$ {\tiny(30)} & $0.889^{*}$ {\tiny(40)} & $0.996^{*}$ {\tiny(3)} & $0.846$ {\tiny(15)} \\
5 & dependence (PS, additive) & 3 & $0.999$ {\tiny(12)} & $0.999^{*}$ {\tiny(8)} & $0.953$ {\tiny(15)} & $0.998^{*}$ {\tiny(10)} & $0.955$ {\tiny(15)} \\
6 & dependence ($\rho(\bs)$, repl.\ $C$) & 3 & $1.000$ {\tiny(10)} & $0.999$ {\tiny(10)} & $0.861^{*}$ {\tiny(20)} & $0.997^{*}$ {\tiny(12)} & $0.941^{*}$ {\tiny(10)} \\
7 & dependence ($\rho(\bs)$, additive) & 3 & $1.000$ {\tiny(25)} & $1.000$ {\tiny(10)} & $0.988$ {\tiny(20)} & $0.995$ {\tiny(25)} & $0.926$ {\tiny(15)} \\
8 & dependence (Fuentes, repl.\ $C$) & 3 & $1.000$ {\tiny(10)} & $1.000$ {\tiny(3)} & $0.848$ {\tiny(25)} & $0.998$ {\tiny(5)} & $\mathbf{0.721^{*}}$ {\tiny(10)} \\
9 & tails ($g$-and-$h$) & 3 & $0.999^{*}$ {\tiny(5)} & $1.000$ {\tiny(10)} & $0.887$ {\tiny(5)} & $0.986$ {\tiny(25)} & $0.942$ {\tiny(5)} \\
10 & skewness ($\lambda(\bs)$) & 3 & $0.994$ {\tiny(25)} & $1.000$ {\tiny(3)} & $0.975$ {\tiny(5)} & $0.998$ {\tiny(25)} & $0.938^{*}$ {\tiny(5)} \\
11 & mean (2 bumps, sd 11) & 3 & $0.976^{*}$ {\tiny(50)} & $0.965^{*}$ {\tiny(50)} & $\mathbf{0.199^{*}}$ {\tiny(20)} & $0.961^{*}$ {\tiny(50)} & $\mathbf{0.267^{*}}$ {\tiny(20)} \\
12 & mean (6 bumps, sd 11) & 3 & $0.960^{*}$ {\tiny(50)} & $0.953^{*}$ {\tiny(50)} & $\mathbf{0.421^{*}}$ {\tiny(40)} & $0.952^{*}$ {\tiny(50)} & $0.846$ {\tiny(15)} \\
\bottomrule
\end{tabular}

\end{table}

\begin{table}[ht]
\centering\small
\caption{\textbf{Brier score} (mean over the 11 exceedance levels), best over
$K>0$ relative to own $K=0$. Bold: $\bar R < 0.75$.}
\label{tab:brier}
\setlength{\tabcolsep}{4pt}
\begin{tabular}{clcccccc}
\toprule
ID & non-stationarity & $n$ & Gaussian & Skew-$t$ K1 & Sym-$t$ K1 T & Skew-$t$ K5 & Sym-$t$ K5 T \\
\midrule
1 & mean (2 bumps, sd 3.7) & 10 & $1.000$ {\tiny(3)} & $0.978^{*}$ {\tiny(50)} & $0.993^{*}$ {\tiny(30)} & $0.984^{*}$ {\tiny(50)} & $0.971$ {\tiny(3)} \\
3 & dependence (low-rank) & 10 & $1.000$ {\tiny(5)} & $0.999$ {\tiny(5)} & $1.000$ {\tiny(3)} & $0.997$ {\tiny(15)} & $0.993$ {\tiny(12)} \\
4 & dependence (PS, repl.\ $C$) & 3 & $0.998$ {\tiny(50)} & $0.998^{*}$ {\tiny(30)} & $0.998^{*}$ {\tiny(30)} & $0.983^{*}$ {\tiny(50)} & $0.940^{*}$ {\tiny(20)} \\
5 & dependence (PS, additive) & 3 & $1.000$ {\tiny(3)} & $0.999$ {\tiny(12)} & $1.000$ {\tiny(8)} & $0.990$ {\tiny(15)} & $0.973$ {\tiny(5)} \\
6 & dependence ($\rho(\bs)$, repl.\ $C$) & 3 & $0.999$ {\tiny(20)} & $0.999^{*}$ {\tiny(5)} & $0.999^{*}$ {\tiny(15)} & $0.983^{*}$ {\tiny(12)} & $0.970$ {\tiny(10)} \\
7 & dependence ($\rho(\bs)$, additive) & 3 & $1.000$ {\tiny(15)} & $0.999$ {\tiny(10)} & $0.998$ {\tiny(15)} & $0.990$ {\tiny(5)} & $0.968$ {\tiny(15)} \\
8 & dependence (Fuentes, repl.\ $C$) & 3 & $1.000^{*}$ {\tiny(3)} & $1.000$ {\tiny(3)} & $0.999$ {\tiny(3)} & $0.997$ {\tiny(12)} & $0.962$ {\tiny(12)} \\
9 & tails ($g$-and-$h$) & 3 & $1.000$ {\tiny(12)} & $1.000$ {\tiny(15)} & $1.000$ {\tiny(3)} & $0.996$ {\tiny(10)} & $0.967$ {\tiny(12)} \\
10 & skewness ($\lambda(\bs)$) & 3 & $0.995$ {\tiny(25)} & $0.999$ {\tiny(3)} & $0.957^{*}$ {\tiny(25)} & $0.994$ {\tiny(25)} & $0.988$ {\tiny(20)} \\
11 & mean (2 bumps, sd 11) & 3 & $1.000$ {\tiny(3)} & $1.000$ {\tiny(3)} & $0.952^{*}$ {\tiny(30)} & $1.004$ {\tiny(3)} & $0.997$ {\tiny(3)} \\
12 & mean (6 bumps, sd 11) & 3 & $0.930$ {\tiny(50)} & $0.925^{*}$ {\tiny(50)} & $0.880^{*}$ {\tiny(50)} & $0.913^{*}$ {\tiny(50)} & $0.884^{*}$ {\tiny(10)} \\
\bottomrule
\end{tabular}

\end{table}

\begin{table}[ht]
\centering\small
\caption{\textbf{Variogram score} ($p = 1/2$), best over $K>0$ relative to own
$K=0$. Bold: $\bar R < 0.75$.}
\label{tab:vario}
\setlength{\tabcolsep}{3.7pt}
\begin{tabular}{clcccccc}
\toprule
ID & non-stationarity & $n$ & Gaussian & Skew-$t$ K1 & Sym-$t$ K1 T & Skew-$t$ K5 & Sym-$t$ K5 T \\
\midrule
1 & mean (2 bumps, sd 3.7) & 10 & $0.994^{*}$ {\tiny(40)} & $0.975^{*}$ {\tiny(40)} & $\mathbf{0.617^{*}}$ {\tiny(15)} & $0.981^{*}$ {\tiny(50)} & $\mathbf{0.712}$ {\tiny(12)} \\
3 & dependence (low-rank) & 10 & $1.000$ {\tiny(5)} & $1.000$ {\tiny(3)} & $0.984$ {\tiny(5)} & $0.999$ {\tiny(25)} & $0.894$ {\tiny(3)} \\
4 & dependence (PS, repl.\ $C$) & 3 & $0.999^{*}$ {\tiny(10)} & $1.000$ {\tiny(15)} & $0.988$ {\tiny(15)} & $0.997$ {\tiny(10)} & $0.990$ {\tiny(3)} \\
5 & dependence (PS, additive) & 3 & $0.999^{*}$ {\tiny(12)} & $0.999^{*}$ {\tiny(8)} & $0.963$ {\tiny(50)} & $0.998$ {\tiny(15)} & $0.986$ {\tiny(15)} \\
6 & dependence ($\rho(\bs)$, repl.\ $C$) & 3 & $0.999$ {\tiny(10)} & $0.999$ {\tiny(50)} & $0.920^{*}$ {\tiny(40)} & $0.997$ {\tiny(12)} & $0.984$ {\tiny(10)} \\
7 & dependence ($\rho(\bs)$, additive) & 3 & $1.000$ {\tiny(10)} & $1.000$ {\tiny(30)} & $0.995$ {\tiny(3)} & $0.995$ {\tiny(5)} & $0.960$ {\tiny(15)} \\
8 & dependence (Fuentes, repl.\ $C$) & 3 & $1.000$ {\tiny(8)} & $1.000$ {\tiny(3)} & $0.887$ {\tiny(40)} & $1.000$ {\tiny(3)} & $0.847$ {\tiny(5)} \\
9 & tails ($g$-and-$h$) & 3 & $1.000$ {\tiny(3)} & $1.000$ {\tiny(5)} & $0.865$ {\tiny(5)} & $0.982$ {\tiny(25)} & $0.947$ {\tiny(5)} \\
10 & skewness ($\lambda(\bs)$) & 3 & $0.990$ {\tiny(40)} & $0.999$ {\tiny(30)} & $1.000$ {\tiny(3)} & $0.993$ {\tiny(30)} & $0.905^{*}$ {\tiny(10)} \\
11 & mean (2 bumps, sd 11) & 3 & $0.963^{*}$ {\tiny(50)} & $0.943^{*}$ {\tiny(50)} & $\mathbf{0.201^{*}}$ {\tiny(10)} & $0.935^{*}$ {\tiny(50)} & $\mathbf{0.244^{*}}$ {\tiny(20)} \\
12 & mean (6 bumps, sd 11) & 3 & $0.938^{*}$ {\tiny(50)} & $0.911^{*}$ {\tiny(50)} & $\mathbf{0.354^{*}}$ {\tiny(40)} & $0.915^{*}$ {\tiny(50)} & $0.909$ {\tiny(15)} \\
\midrule
\multicolumn{8}{l}{\emph{Non-$t$ extension (Gaussian-GP errors; skew-$t$ is misspecified)}} \\
13 & mean, Franke (multiscale) & 3 & $0.912^{*}$ {\tiny(50)} & $0.913^{*}$ {\tiny(50)} & $1.164^{*}$ {\tiny(30)} & $0.912^{*}$ {\tiny(50)} & $1.085$ {\tiny(30)} \\
14 & mean, curved ridge & 3 & $0.891^{*}$ {\tiny(40)} & $0.891^{*}$ {\tiny(40)} & $2.092^{*}$ {\tiny(50)} & $0.913^{*}$ {\tiny(40)} & $2.571^{*}$ {\tiny(40)} \\
15 & mean, annulus (Ricker) & 3 & $1.032^{*}$ {\tiny(50)} & $1.031^{*}$ {\tiny(50)} & $1.448^{*}$ {\tiny(30)} & $1.035^{*}$ {\tiny(50)} & $1.444^{*}$ {\tiny(40)} \\
16 & mean, sigmoid front & 10 & $0.941^{*}$ {\tiny(20)} & $0.941^{*}$ {\tiny(50)} & $0.998$ {\tiny(3)} & $0.938^{*}$ {\tiny(20)} & $1.000$ {\tiny(3)} \\
17 & mean, GP draw & 3 & $1.000$ {\tiny(50)} & $0.998$ {\tiny(50)} & $1.205$ {\tiny(40)} & $0.983$ {\tiny(50)} & $1.218$ {\tiny(40)} \\
18 & mean, Franke + sinh-arcsinh err. & 3 & $0.905^{*}$ {\tiny(50)} & $0.906^{*}$ {\tiny(50)} & $\mathbf{0.726^{*}}$ {\tiny(30)} & $0.919^{*}$ {\tiny(50)} & $0.854$ {\tiny(30)} \\
19 & mean, poly2 (bowl+saddle) & 3 & $0.935^{*}$ {\tiny(30)} & $0.935^{*}$ {\tiny(50)} & $1.542$ {\tiny(40)} & $0.928^{*}$ {\tiny(50)} & $1.020$ {\tiny(40)} \\
20 & mean, transform mix & 30 & $0.955^{*}$ {\tiny(50)} & $0.953^{*}$ {\tiny(50)} & $\mathbf{0.389^{*}}$ {\tiny(15)} & $\mathbf{0.631^{*}}$ {\tiny(50)} & $\mathbf{0.402^{*}}$ {\tiny(50)} \\
21 & mean, non-smooth (hinges) & 3 & $0.934^{*}$ {\tiny(30)} & $0.932^{*}$ {\tiny(30)} & $1.055$ {\tiny(30)} & $0.904^{*}$ {\tiny(30)} & $1.209$ {\tiny(30)} \\
\bottomrule
\end{tabular}

\end{table}

\begin{table}[p]
\centering\footnotesize
\setlength{\tabcolsep}{3.5pt}
\caption{\textbf{Relative Brier score at every $\Kmrts$}, K1-knot methods (Gaussian, Skew-$t$ K1, Sym-$t$ K1 T), skew-$t$ settings (vs the same method's own $K=0$, mean over datasets; $^{*}$ marks $|\bar R - 1| > 2\,\widehat{\mathrm{SE}}$). A dash means that $K$ was not fitted for that setting (settings 13--21 use only $K \in \{30,40,50\}$). Isolated large values of Sym-$t$ K5 T are the exponential-fallback outliers discussed in the Caveats and are reported as-is.}
\label{tab:brierK}
\begin{tabular}{ccccccccccccc}
\toprule
ID & $n$ & $K{=}3$ & $K{=}5$ & $K{=}8$ & $K{=}10$ & $K{=}12$ & $K{=}15$ & $K{=}20$ & $K{=}25$ & $K{=}30$ & $K{=}40$ & $K{=}50$ \\
\midrule
\multicolumn{13}{l}{\emph{Gaussian}} \\[1pt]
1 & 10 & $1.000$ & $1.016^{*}$ & $1.023^{*}$ & $1.022^{*}$ & $1.019^{*}$ & $1.016^{*}$ & $1.009^{*}$ & $1.008^{*}$ & $1.005$ & $1.003$ & $1.003$ \\
3 & 10 & $1.000$ & $1.000$ & $1.001^{*}$ & $1.001^{*}$ & $1.000$ & $1.000$ & $1.000$ & $1.001$ & $1.001$ & $1.001^{*}$ & $1.002^{*}$ \\
4 & 3 & $1.000$ & $0.999^{*}$ & $0.999^{*}$ & $1.000$ & $1.000$ & $1.000$ & $0.999$ & $0.999$ & $0.999$ & $0.999$ & $0.998$ \\
5 & 3 & $1.000$ & $1.001$ & $1.001$ & $1.000$ & $1.001^{*}$ & $1.000$ & $1.001^{*}$ & $1.002^{*}$ & $1.002^{*}$ & $1.001^{*}$ & $1.004^{*}$ \\
6 & 3 & $1.000$ & $1.000$ & $1.000$ & $1.000$ & $1.000$ & $1.000$ & $0.999$ & $1.000$ & $0.999$ & $0.998^{*}$ & $0.998$ \\
7 & 3 & $1.000$ & $1.000$ & $1.000$ & $1.000$ & $1.000$ & $1.000$ & $1.002^{*}$ & $1.001$ & $1.000$ & $1.000$ & $1.000$ \\
8 & 3 & $1.000^{*}$ & $1.001$ & $1.001$ & $1.001$ & $1.000$ & $1.000$ & $1.002$ & $1.002$ & $1.002$ & $1.003$ & $1.004^{*}$ \\
9 & 3 & $1.000^{*}$ & $1.000$ & $1.000$ & $1.000^{*}$ & $1.000$ & $1.000$ & $1.000$ & $1.001$ & $1.001$ & $1.002$ & $1.001$ \\
10 & 3 & $1.000$ & $1.003$ & $1.008$ & $1.004$ & $1.000$ & $0.997$ & $0.995$ & $0.995$ & $0.989$ & $0.988$ & $0.988$ \\
11 & 3 & $1.000$ & $1.059^{*}$ & $1.078^{*}$ & $1.101^{*}$ & $1.114^{*}$ & $1.095^{*}$ & $1.069^{*}$ & $1.051^{*}$ & $1.022$ & $1.001$ & $1.000$ \\
12 & 3 & $1.000$ & $1.003^{*}$ & $1.000$ & $0.993$ & $0.980$ & $1.019$ & $1.049$ & $1.045$ & $0.988$ & $0.947$ & $0.930$ \\
\midrule
\multicolumn{13}{l}{\emph{Skew-$t$ K1}} \\[1pt]
1 & 10 & $1.000$ & $0.996$ & $1.006$ & $1.006$ & $1.001$ & $0.996$ & $0.987^{*}$ & $0.985^{*}$ & $0.981^{*}$ & $0.978^{*}$ & $0.978^{*}$ \\
3 & 10 & $1.000$ & $0.999$ & $1.000$ & $1.000$ & $1.000$ & $1.000$ & $1.000$ & $1.000$ & $1.000$ & $1.000$ & $1.000$ \\
4 & 3 & $1.000^{*}$ & $1.000$ & $0.999$ & $1.000$ & $1.000$ & $1.000$ & $0.999$ & $0.999$ & $0.998^{*}$ & $0.999$ & $0.999$ \\
5 & 3 & $1.000$ & $1.000$ & $0.999$ & $1.000$ & $0.999$ & $1.000$ & $1.000$ & $1.001$ & $1.001$ & $1.001$ & $1.002$ \\
6 & 3 & $1.000$ & $0.999^{*}$ & $0.999^{*}$ & $1.000$ & $1.001^{*}$ & $0.999$ & $1.000$ & $1.000$ & $1.000$ & $0.999$ & $1.000$ \\
7 & 3 & $1.000$ & $1.000$ & $1.001$ & $0.999$ & $1.000$ & $1.000$ & $1.001$ & $1.000$ & $1.000$ & $1.001$ & $1.001$ \\
8 & 3 & $1.000$ & $1.000$ & $1.001$ & $1.000$ & $1.001$ & $1.001^{*}$ & $1.001$ & $1.000$ & $1.001$ & $1.001^{*}$ & $1.002^{*}$ \\
9 & 3 & $1.001$ & $1.001$ & $1.001^{*}$ & $1.001$ & $1.001$ & $1.000$ & $1.001$ & $1.001$ & $1.001$ & $1.001$ & $1.001$ \\
10 & 3 & $0.999$ & $1.003$ & $1.008^{*}$ & $1.008$ & $1.005$ & $1.003$ & $1.002$ & $1.001$ & $0.999$ & $0.998$ & $0.998$ \\
11 & 3 & $1.000$ & $1.052^{*}$ & $1.074^{*}$ & $1.103^{*}$ & $1.163^{*}$ & $1.138^{*}$ & $1.110^{*}$ & $1.086^{*}$ & $1.048$ & $1.017$ & $1.014$ \\
12 & 3 & $1.000$ & $1.005^{*}$ & $0.999$ & $0.992$ & $0.965$ & $1.022$ & $1.024$ & $1.020$ & $0.980$ & $0.946$ & $0.925^{*}$ \\
\midrule
\multicolumn{13}{l}{\emph{Sym-$t$ K1 T}} \\[1pt]
1 & 10 & $1.000$ & $1.003$ & $1.006$ & $1.008^{*}$ & $1.003$ & $1.001$ & $0.996$ & $0.994^{*}$ & $0.993^{*}$ & $0.993^{*}$ & $0.994^{*}$ \\
3 & 10 & $1.000$ & $1.000$ & $1.001$ & $1.001$ & $1.001$ & $1.001$ & $1.001$ & $1.001$ & $1.002$ & $1.002$ & $1.004^{*}$ \\
4 & 3 & $1.000$ & $0.998^{*}$ & $0.999$ & $0.999$ & $1.000$ & $0.999$ & $1.000$ & $0.999$ & $0.998^{*}$ & $1.000$ & $0.999$ \\
5 & 3 & $1.000^{*}$ & $1.000$ & $1.000$ & $1.000$ & $1.000$ & $1.001$ & $1.000$ & $1.001$ & $1.002^{*}$ & $1.004$ & $1.007$ \\
6 & 3 & $1.000$ & $1.002^{*}$ & $1.000$ & $1.000$ & $1.001$ & $0.999^{*}$ & $0.999$ & $1.000$ & $1.001$ & $0.997^{*}$ & $1.000$ \\
7 & 3 & $1.000^{*}$ & $1.000$ & $1.000$ & $1.000$ & $1.000$ & $0.998$ & $1.001$ & $1.000$ & $1.000$ & $0.999$ & $1.001$ \\
8 & 3 & $0.999$ & $0.999$ & $0.999$ & $1.002$ & $1.000$ & $1.001$ & $1.003$ & $1.008$ & $1.009$ & $1.012$ & $1.010^{*}$ \\
9 & 3 & $1.000$ & $1.001$ & $1.001$ & $1.000$ & $1.001$ & $1.001^{*}$ & $1.000$ & $1.001^{*}$ & $1.002$ & $0.998$ & $1.000$ \\
10 & 3 & $0.999$ & $0.993^{*}$ & $0.980^{*}$ & $0.985$ & $0.975$ & $0.970$ & $0.963^{*}$ & $0.957^{*}$ & $0.949^{*}$ & $0.954^{*}$ & $0.949^{*}$ \\
11 & 3 & $0.999$ & $1.016$ & $1.021$ & $1.039^{*}$ & $1.008$ & $0.982$ & $0.977^{*}$ & $0.959^{*}$ & $0.952^{*}$ & $0.954^{*}$ & $0.959^{*}$ \\
12 & 3 & $1.000$ & $1.006^{*}$ & $0.997$ & $0.907^{*}$ & $0.893^{*}$ & $0.899^{*}$ & $0.896^{*}$ & $0.918^{*}$ & $0.892^{*}$ & $0.896^{*}$ & $0.880^{*}$ \\
\bottomrule
\end{tabular}
\end{table}


\subsection{Finding 1: gains appear only where the truth has mean structure ---
and this is provable a priori}

The only large entries in any table are settings 1, 11 and 12 (mean
structure), methods 3 and 5: energy and variogram reductions of 34--42\%
(setting 1), 73--80\% (setting 11) and 58\% (setting 12). Every dependence,
tails and skewness setting is flat to within a few percent, in \emph{all
three} scores.

This dichotomy can be established without MCMC. Because the MRTS coefficient
$\beta$ is pooled over time, the most any mean basis can deliver is the
time-averaged true mean $\bar\mu(\bs)$. Define the oracle relative Brier
\begin{equation}
R^\star \;=\; \frac{\mathrm{BS}\big(\bar\mu\big)}{\mathrm{BS}\big(\text{spatial
constant}\big)},
\end{equation}
both evaluated under the \emph{true} skew-$t$ predictive law. $R^\star$ is a
floor: no $K$, sample size or sampler can push the relative Brier below it.
Table~\ref{tab:oracle} shows $R^\star \approx 1$ for nine of eleven settings:
their Brier nulls are design properties, not evidence about MRTS. Setting 10 is
the striking case: $R^\star = 1.12 > 1$ --- supplying the \emph{true} mean makes
the Brier score \emph{worse}, because the mean shift induced by
$\lambda(\bs)$ is a symptom of a skewness change that the constant-$\lambda$
model cannot represent; correcting the symptom while the disease remains
actively misleads the tail probabilities.

\begin{table}[ht]
\centering\small
\caption{Oracle Brier floor $R^\star$ per setting
(\texttt{oracle\_brier\_ceiling.R}; Monte Carlo, 3 datasets).}
\label{tab:oracle}
\begin{tabular}{lccc}
\toprule
setting & $\mathrm{sd}(\bar\mu)$ & $R^\star$ & verdict \\
\midrule
1 invariant          & 3.68  & 0.856 & weak headroom (14\%) \\
2 varying            & 0.42  & 1.000 & dead end \\
3 ns\_dependence     & 0.41  & 0.999 & dead end \\
4 ps\_cov            & 0.00  & 1.000 & dead end \\
5 ps\_add            & 0.39  & 0.999 & dead end \\
6 covdep\_cov        & 0.00  & 1.000 & dead end \\
7 covdep\_add        & 0.41  & 1.000 & dead end \\
8 fuentes\_cov       & 0.00  & 1.000 & dead end \\
9 gh\_marginal       & 0.43  & 1.000 & dead end \\
10 lambda\_varying   & 2.80  & \textbf{1.120} & \textbf{harmful} \\
11 invariant\_strong & 11.04 & 0.439 & positive control \\
12 mrts\_multibump   & 11.00 & 0.495 & positive control \\
\bottomrule
\end{tabular}
\end{table}

\subsection{Finding 2: the gains are carried entirely by the thresholded
methods}

In Table~\ref{tab:energy}, setting 11 (mean sd $=11$), the two POT-thresholded
symmetric-$t$ methods reduce the energy score by 73--80\%, while the Gaussian
and both skew-$t$ methods stay within 2\% of their $K=0$ baselines --- despite
those same skew-$t$ fits recovering the mean surface almost exactly
(recovery RMSE $14.0 \to 0.66$ for method 2 at $K=25$) and recovering every DGP
parameter (Section~\ref{sec:anomaly}). Table~\ref{tab:perK} shows the full $K$
profile.

\begin{table}[ht]
\centering\small
\caption{Setting 11, relative score vs $\Kmrts$ ($n=3$): the thresholded
symmetric-$t$ methods against the unthresholded skew-$t$.}
\label{tab:perK}
\footnotesize
\setlength{\tabcolsep}{3.5pt}
\begin{tabular}{llccccccccccc}
\toprule
score & method & $K{=}3$ & $K{=}5$ & $K{=}8$ & $K{=}10$ & $K{=}12$ & $K{=}15$ & $K{=}20$ & $K{=}25$ & $K{=}30$ & $K{=}40$ & $K{=}50$ \\
\midrule
\multirow{3}{*}{energy}
 & Skew-$t$ K1 (no thr.)   & $0.999$ & $1.013$ & $1.022$ & $1.027$ & $1.023$ & $1.017$ & $0.996$ & $0.984$ & $0.978$ & $0.966$ & $0.965$ \\
 & Sym-$t$ K1 (thr.)       & $0.990$ & $\mathbf{0.263}$ & $0.266$ & $0.229$ & $0.201$ & $0.244$ & $\mathbf{0.199}$ & $0.256$ & $0.233$ & $0.268$ & $0.305$ \\
 & Sym-$t$ K5 (thr.)       & $1.001$ & $0.317$ & $0.539$ & $0.374$ & $0.342$ & $0.462$ & $\mathbf{0.267}$ & $0.333$ & $0.359$ & $0.906$ & $0.860$ \\
\midrule
\multirow{3}{*}{Brier}
 & Skew-$t$ K1 (no thr.)   & $1.000$ & $1.052$ & $1.074$ & $1.103$ & $\mathbf{1.163}$ & $1.138$ & $1.110$ & $1.086$ & $1.048$ & $1.017$ & $1.014$ \\
 & Sym-$t$ K1 (thr.)       & $0.999$ & $1.016$ & $1.021$ & $1.039$ & $1.008$ & $0.982$ & $0.977$ & $0.959$ & $0.952$ & $0.954$ & $0.959$ \\
 & Sym-$t$ K5 (thr.)       & $0.997$ & $1.078$ & $1.182$ & $1.114$ & $1.163$ & $1.120$ & $1.104$ & $1.133$ & $1.030$ & $1.086$ & $1.063$ \\
\bottomrule
\end{tabular}
\end{table}

Methods 3/5 are simultaneously (a) POT-thresholded and (b) symmetric-$t$, so a
disambiguating pair was added: methods 11/12, \emph{identical to methods 3/5
except} $\text{threshold} = 0$, fitted on setting 1 with $n = 10$
(Table~\ref{tab:disambig}). The verdict is unambiguous: the unthresholded
symmetric-$t$ methods are \emph{flat} (energy ratio $0.99$ at every $K$,
indistinguishable from the skew-$t$ reference), while their thresholded twins
gain 35--40\%. \textbf{The threshold, not the symmetry, is the driver.}

\begin{table}[ht]
\centering\small
\caption{Disambiguation on setting 1 ($n = 10$): relative energy score vs
$\Kmrts$, and the \emph{absolute} energy score at $K = 0$. Method pairs
(3,\,11) and (5,\,12) differ \emph{only} in the POT threshold.}
\label{tab:disambig}
\begin{tabular}{llccccc|c}
\toprule
 & method & $K{=}5$ & $K{=}10$ & $K{=}15$ & $K{=}20$ & $K{=}25$ & abs.\ ES at $K{=}0$ \\
\midrule
3  & Sym-$t$ K1, thr.\ $q(0.80)$ & $0.630$ & $0.617$ & $0.589$ & $\mathbf{0.579}$ & $0.596$ & $21.8$ \\
11 & Sym-$t$ K1, no threshold    & $0.996$ & $0.997$ & $0.991$ & $0.990$ & $0.989$ & $8.5$ \\
\midrule
5  & Sym-$t$ K5, thr.\ $q(0.80)$ & $0.688$ & $1.113$ & $0.773$ & $0.718$ & $\mathbf{0.697}$ & $137.1$ \\
12 & Sym-$t$ K5, no threshold    & $0.994$ & $1.008$ & $0.993$ & $0.997$ & $0.997$ & $9.3$ \\
\midrule
2  & Skew-$t$ K1 (reference)     & $0.994$ & $0.994$ & $0.985$ & $0.983$ & $0.982$ & $8.5$ \\
\bottomrule
\end{tabular}
\end{table}

The absolute column reframes the finding. The thresholded baselines are
\emph{catastrophically} worse to begin with: censoring everything below
$q(0.80)$ leaves the constant-mean fit with almost no information about the
spatial level of the field, and its out-of-sample predictive field degrades by
a factor $2.6$ (K1) to $16$ (K5) in energy score. MRTS repairs the mean that
the censored likelihood cannot learn --- a partial \emph{rescue}, not a
synergy: even at the best $K$, method 3 reaches
$21.8 \times 0.579 \approx 12.6$, still well above the $8.5$ that methods 2/11
achieve \emph{without any MRTS at all}. In absolute terms the unthresholded
methods win at every $K$; the correct reading of Finding 2 is therefore
``\emph{only the thresholded methods have room to be helped},'' and the
practical implication is for POT-style threshold models specifically: if the
mean is not modelled well, censored fitting amplifies the damage, and a mean
basis buys back part of it.

\subsection{Finding 3: the Brier score is nearly blind to mean recovery, and
can punish it}\label{sec:anomaly}

Setting 11, method 2 (skew-$t$ K1), dataset 1, posterior means
(truth: $\lambda = 3$, $\E[\sigma] = 2.26$, $\rho = 1$, $\gamma = 0.9$):

\begin{center}\small
\begin{tabular}{lccccc}
\toprule
$\Kmrts$ & recovery RMSE & $\hat\lambda$ & $\E[\hat\sigma]$ & $\hat\rho$ & $\hat\gamma$ \\
\midrule
0  & 14.01 & 0.15 & 11.16 & 2.04 & 0.99 \\
10 & 2.51  & 2.04 & 3.01  & 1.21 & 0.91 \\
25 & \textbf{0.66} & \textbf{3.18} & \textbf{1.94} & \textbf{0.96} & \textbf{0.94} \\
\bottomrule
\end{tabular}
\end{center}

At $K=0$ the model cannot represent the sd-11 mean surface and absorbs the
spatial variance into an inflated $\sigma$ (5$\times$ truth) with $\lambda$
crushed to zero. At $K=25$ \emph{every} parameter is recovered. Yet the $K=25$
Brier score is \emph{worse} than $K=0$ by 8.6\% (Table~\ref{tab:perK}), and the
same reversal holds for methods 1, 4, 5. A correctly specified model loses to a
badly misspecified one on this score. The working hypothesis --- not yet
verified --- is an out-of-sample interaction: the sharpened predictive
distribution ($\hat\sigma$ correct instead of inflated) exposes the residual
error of the MRTS mean \emph{extrapolated} to the held-out sites, whereas the
misspecified $K=0$ fit hides that error inside an accidentally conservative
(inflated-$\sigma$) predictive distribution.

\paragraph{Setting 12 resolves the anomaly.}
The rich-mean setting 12 (six bumps, sd 11) was run out to $\Kmrts = 50$,
where the noise-free MRTS projection error of its surface finally becomes
small ($0.71$ at $K=5$, $0.13$ at $K=25$, $0.057$ at $K=50$).
The Brier profile follows exactly the pattern the extrapolation mechanism
predicts: flat-to-worse at intermediate $K$ (ratios $1.00$--$1.05$ for
$K \le 25$, where the basis is still too coarse and the residual mean error is
exposed by the sharpened predictive), then \emph{improving} once the basis is
fine enough:

\begin{center}\small
\begin{tabular}{lcccc|c}
\toprule
relative Brier, setting 12 & $K{=}15$ & $K{=}25$ & $K{=}30$ & $K{=}50$ &
recovery at $K{=}50$ \\
\midrule
Gaussian        & $1.019$ & $1.045$ & $0.988$ & $\mathbf{0.930}$ & $0.512$ \\
Skew-$t$ K1     & $1.022$ & $1.020$ & $0.980$ & $\mathbf{0.925}$ & $0.073$ \\
Sym-$t$ K1 thr. & $0.899$ & $0.918$ & $0.892$ & $\mathbf{0.880}$ & $0.422$ \\
Skew-$t$ K5     & $0.976$ & $1.022$ & $0.965$ & $\mathbf{0.913}$ & $0.268$ \\
\bottomrule
\end{tabular}
\end{center}

(Method 5 is excluded: its exponential-fallback instability blows up one of the
three datasets at large $K$, ratio $2.7$.) The $K{=}50$ gains of 7--12\% are
real but far above the oracle floor $R^\star = 0.495$ --- the mean channel is
recovered (skew-$t$ K1 recovery ratio $0.073$), yet most of the theoretically
available Brier headroom is spent compensating the sharpness--extrapolation
trade-off. Practically: \emph{for judging mean-structure recovery, the Brier
score remains the wrong instrument --- it responds only when the basis is fine
enough to make the residual extrapolation error negligible; energy/variogram
scores (or the recovery RMSE directly) are the informative ones.}

\paragraph{The large-$K$ limit discriminates the mechanism.}
If the reversal were simply a matter of ``$K$ large enough,'' setting 11 ---
whose 2-bump surface the basis captures by $K \approx 10$ --- should also turn
positive by $K = 50$. It does not (Table~\ref{tab:perK}, rightmost columns):
the mid-$K$ deterioration decays monotonically ($1.16$ at $K{=}12$ to
$1.01$--$1.00$ at $K{=}50$ for the skew-$t$ and Gaussian methods) but never
crosses below parity, leaving the oracle headroom $R^\star = 0.439$ entirely
unused. The weak-signal setting 1 behaves the same way at large $K$ (skew-$t$
K1: $0.985$ at $K{=}25$, $0.978$ at $K{=}50$ --- smooth, small, no regime
change). Same $K$ grid, same machinery, opposite outcome to setting 12: the
Brier reversal requires mean structure that \emph{only} a fine basis can
resolve; extra columns on an already-resolved surface buy nothing and merely
let the sharpness--extrapolation penalty fade.

\section{Caveats}

\begin{itemize}
\item Settings 4--11 are pilots with $n = 3$; settings 1 and 3 have $n = 10$.
Starred cells with $n=3$ should be read as indicative. The small-$n$
\emph{best-over-}$K$ statistic also has a selection bias toward values $<1$;
flat profiles (e.g.\ Sym-$t$ K5 rows in dependence settings, best
$\approx 0.72$--$0.85$ with non-monotone $K$ profiles) are consistent with
noise plus this selection effect.
\item Method 5 (Sym-$t$ K5 T) occasionally triggers its exponential-covariance
fallback and produces isolated outlier fits (e.g.\ setting 4, $K=25$: energy
ratio 15.8 in one dataset). Best-over-$K$ cells are robust to this; per-$K$
curves at $n=3$ are not.
\item All $K$ comparisons are paired (shared seeds across $K$), so
within-method ratios are considerably more precise than the between-method
comparisons.
\item The recovery RMSE (not tabulated here) has a known bimodality on
flat-mean settings caused by the intercept--$\lambda$ ridge; see
\texttt{nonsta\_settings\_design.md}.
\end{itemize}

\section{The setting-12 showcase: why a large $K$ is worth having}

Setting 12's six-bump surface was designed so that the MRTS approximation
improves \emph{gradually} with $K$ (noise-free projection error $0.71$ at
$K{=}5$, $0.19$ at $K{=}15$, $0.057$ at $K{=}50$), in contrast to the 2-bump
settings which the basis captures by $K{=}5$. The fitted results track that
curve: skew-$t$ K1 recovery ratios are $0.93\ (K{=}5) \to 0.49\ (K{=}10) \to
0.14\ (K{=}25) \to 0.073\ (K{=}50)$ --- the extended elbow, realised in MCMC.
Two by-products: (a) the Brier anomaly resolution above; (b) a resolution
ceiling --- a higher-frequency mean component (3 cycles per axis) is
\emph{unrecoverable} on the 100-train/44-test geometry (its projection error
\emph{rises} with $K$), i.e.\ MRTS's usable resolution is capped by site
density, and the six-bump surface sits just inside that cap.

\section{In progress}

\begin{enumerate}
\item \textbf{Non-$t$ extension (settings 13--21)}: all settings so far share
the skew-$t$ core, so at large $K$ the fitted skew-$t$ is \emph{correctly
specified} --- a confound for every ``MRTS helps/hurts'' conclusion. The
extension drops the skew-$t$ generator entirely: nine nonlinear mean surfaces
(Franke, curved ridge, annulus, sigmoid front, GP draw, polynomial, transform
mixtures, non-smooth hinges; orthogonalised against the $K=0$ design and
standardised to sd 11) with Gaussian-GP errors, calibrated so that
$R^\star \approx 0.45$ throughout. Data generation in
\texttt{setup\_nonsta.R} $\to$ \texttt{simdata\_nonsta\_ext.RData}.
\item \textbf{$n$ upgrade}: settings 4--12 are $n=3$ pilots; the headline
contrasts (11, 12, and the disambiguation) merit $n = 10$.
\end{enumerate}

\end{document}
